%----------------------------------------------------------------------
\appendix
\section{Coefficient comparison with the Hickerson--Mortenson theta term}
\label{app:HM-coeff-comparison}
%----------------------------------------------------------------------

We keep the conventions of the paper:
\[
  C_2(z)=\binom z2=\frac{z(z-1)}2,\qquad
  B=n+p,\qquad
  \Delta=p(2n+p),\qquad
  N=n\Delta,
\]
\[
  A_0=\frac{np(2n+p+1)}2,\qquad
  H(R,S)=nC_2(R)+BRS+nC_2(S),\qquad
  \sgtwo(u,v)=\frac{\sg(u)+\sg(v)}2,
\]
and
\[
  j(z;q)=\sum_{m\in\mathbb Z}(-1)^m q^{C_2(m)}z^m,\qquad
  \sg(0)=1.
\]
Let
\[
  \alpha=\frac{n-1}{2},\qquad \beta=\frac{n+1}{2},\qquad
  \varepsilon=\{\alpha\}\in\{0,1/2\}.
\]
For \(r^*,s^*\in\{0,\ldots,p-1\}\), put
\[
  \widetilde r=r^*+\varepsilon,\qquad
  \widetilde s=s^*+\varepsilon,\qquad
  \rho=\widetilde r-\alpha,\qquad
  \sigma=\widetilde s+\beta.
\]
Then \(\rho,\sigma\in\mathbb Z\), and
\(\widetilde r,\widetilde s\in \varepsilon+\{0,\ldots,p-1\}\subset[0,p)\).
Also
\[
  \widetilde s-\widetilde r=\sigma-\rho-n.
\]
The Hickerson--Mortenson theta correction, specialized to
\((z_0,z_1)=(-1,-1)\), can be rewritten as
\[
\begin{aligned}
\theta^{\rm HM}_{n,p}
=&\sum_{r^*=0}^{p-1}\sum_{s^*=0}^{p-1}
 q^{H(\rho,\sigma)}(-x)^\rho(-y)^\sigma
 j\!\left(-q^{np(\sigma-\rho-n)}\frac{x^n}{y^n};q^{np^2}\right) \\
&\qquad\qquad\times
 \frac{J_Q^3\,j(d_1d_2;Q)}{j(d_1;Q)j(d_2;Q)},
\end{aligned}
\tag{A.1}
\]
where
\[
  Q=q^{p\Delta}=q^{p^2(2n+p)},
\]
\[
  d_1=q^{\Delta\rho+A_0-np}\frac{(-y)^B}{(-x)^n},\qquad
  d_2=q^{\Delta\sigma-A_0}\frac{(-x)^B}{(-y)^n}.
\]
Indeed,
\[
  d_1d_2=q^{\Delta(\rho+\sigma)-np}(-x)^p(-y)^p,
\]
which is the second numerator theta argument in HM's formula, since
\[
  \Delta(\widetilde r+\widetilde s)+pB
  =\Delta(\rho+\sigma-1)+pB
  =\Delta(\rho+\sigma)-np.
\]

\begin{lemma}[Kronecker partial fraction]\label{lem:kronecker}
With $j(z;Q)=\sum_m(-1)^mQ^{\binom{m}{2}}z^m$,
\[
  \frac{J_Q^3\,j(ab;Q)}{j(a;Q)\,j(b;Q)}
  =\sum_{\eta\in\mathbb Z}\frac{b^\eta}{1-aQ^\eta}.
\tag{A.2}
\]
\end{lemma}

\begin{proof}
We give the coefficientwise proof used in the Lean development.  Put
\[
  \operatorname{PF}(a)=
  \sum_{k\in\mathbb Z}
    \frac{(-1)^kQ^{\binom{k+1}{2}}}{1-aQ^k},
\]
with every denominator expanded in its signed $Q$-first Hahn branch.
The coefficientwise partial-fraction theorem
\[
  j(a;Q)\operatorname{PF}(a)=J_Q^3
\]
follows by a geometric telescope: all nonconstant $a$-coefficients cancel,
while the constant coefficient is Jacobi's cube identity.

Let $K(a,b)=\sum_{\eta\in\mathbb Z}b^\eta/(1-aQ^\eta)$.
The coefficient of $b^B$ in $j(b;Q)K(a,b)$ is
\[
  \sum_{s\in\mathbb Z}
    \frac{(-1)^sQ^{\binom{s}{2}}}{1-aQ^{B-s}}
  =\operatorname{PF}(aQ^B),
\]
where the second equality uses $k=-s$ and
$\binom{s}{2}=\binom{k+1}{2}$.  The inverse-coefficient telescope also gives
\[
  \operatorname{PF}(aQ^B)
  =(-1)^BQ^{\binom{B}{2}}a^B\operatorname{PF}(a)
  \qquad(B\in\mathbb Z).
\]
In the one-step relation, the only boundary term is a coefficient of
$j(1;Q)$, whose two roots $k$ and $-k-1$ cancel.  The last display is exactly
the coefficient of $b^B$ in $j(ab;Q)\operatorname{PF}(a)$.  Hence
\[
  j(b;Q)K(a,b)=j(ab;Q)\operatorname{PF}(a)
\]
coefficientwise.  Multiplication by $j(a;Q)$ and the partial-fraction theorem
above gives
\[
  j(a;Q)j(b;Q)K(a,b)=J_Q^3j(ab;Q);
\]
cancellation in the
corresponding Hahn fraction field proves (A.2).  This is the classical
Kronecker identity (cf.~Weil~\cite{oWe76}, Chapter~IV,
and~\cite{jHiMo14}, Proposition~5.7 and Lemma~5.10), but no residue or
Liouville argument is needed here.
\end{proof}
Applying Lemma~\ref{lem:kronecker} with $a=d_1$, $b=d_2$ gives
\[
  \frac{J_Q^3j(d_1d_2;Q)}{j(d_1;Q)j(d_2;Q)}
  =\sum_{\eta\in\mathbb Z}\frac{d_2^\eta}{1-Q^\eta d_1}.
\tag{A.3}
\]
Next expand the first numerator theta factor by the Jacobi triple product:
\[
\begin{aligned}
 j\!\left(-q^{np(\sigma-\rho-n)}\frac{x^n}{y^n};q^{np^2}\right)
 &=\sum_{\iota\in\mathbb Z}
 (-1)^\iota q^{np^2 C_2(\iota)}
 \left(-q^{np(\sigma-\rho-n)}\frac{x^n}{y^n}\right)^\iota \\
 &=\sum_{\iota\in\mathbb Z}
 q^{np^2 C_2(\iota)+np(\sigma-\rho-n)\iota}
 (-x)^{n\iota}(-y)^{-n\iota}.
\end{aligned}
\tag{A.4}
\]
Here the signs cancel because \((-1)^\iota(-1)^\iota=1\).

For the remaining denominator, set
\[
  M_{\rho,\eta}=Q^\eta d_1
  =q^{p\Delta\eta+\Delta\rho+A_0-np}\frac{(-y)^B}{(-x)^n}.
\]
Its leading \(q\)-exponent is
\[
  \Gamma_{\rho,\eta}=p\Delta\eta+\Delta\rho+A_0-np.
\]
\emph{Formal--analytic bridge.}
On the unit torus $|x|=|y|=1$ with $|q|<1$,
$|M_{\rho,\eta}|=|q|^{\Gamma_{\rho,\eta}}$, so
$|M_{\rho,\eta}|<1$ iff $\Gamma_{\rho,\eta}>0$.
Hence the Hahn branch of $(1-M_{\rho,\eta})^{-1}$
coincides with the convergent geometric series.  Thus the coefficientwise
proof of Lemma~\ref{lem:kronecker} specializes to the analytic quotient on
this torus, with the same Hahn branch convention.
The signed Hahn geometric expansion gives
\[
  (1-M_{\rho,\eta})^{-1}
  =\sum_{\ell\in\mathbb Z} c_{\rho,\eta}(\ell)M_{\rho,\eta}^{\ell},\qquad
  c_{\rho,\eta}(\ell)=\frac{\sg(\ell)+\sg(\Gamma_{\rho,\eta})}{2}.
\tag{A.5}
\]

Combining (A.1)--(A.5), a term indexed by
\((\rho,\sigma,\iota,\eta,\ell)\) has monomial part
\[
  (-x)^{\rho+n\iota+B\eta-n\ell}
  (-y)^{\sigma-n\iota-n\eta+B\ell}.
\]
Thus the coefficient of \((-x)^R(-y)^S\) is constrained by
\[
  R=\rho+n\iota+B\eta-n\ell,\qquad
  S=\sigma-n\iota-n\eta+B\ell.
\tag{A.6}
\]
Define
\[
  w=\ell-\iota-\eta.
\tag{A.7}
\]
Then (A.6) implies
\[
  U_w:=R+nw+\alpha=\rho+p\eta+\alpha=\widetilde r+p\eta,
\tag{A.8}
\]
\[
  V_w:=S-nw-\beta=\sigma-\beta+p\ell=\widetilde s+p\ell.
\tag{A.9}
\]
Conversely, for fixed \(R,S,w\), equations (A.8)--(A.9) determine the
contributing residue pair and summation indices uniquely.  Indeed, since
\(\widetilde r\in\varepsilon+\{0,\ldots,p-1\}\), there is a unique
\(\widetilde r\) in that residue interval and a unique \(\eta\in\mathbb Z\)
such that
\[
  R+nw+\alpha=\widetilde r+p\eta.
\]
Similarly, there is a unique \(\widetilde s\in\varepsilon+\{0,\ldots,p-1\}\)
and a unique \(\ell\in\mathbb Z\) such that
\[
  S-nw-\beta=\widetilde s+p\ell.
\]
Then \(\rho=\widetilde r-\alpha\), \(\sigma=\widetilde s+\beta\), and
(A.7) determines the unique remaining index
\[
  \iota=\ell-\eta-w.
\]
Thus coefficient extraction over the HM residue variables and the bilateral
indices is equivalent to summing over \(w\in\mathbb Z\).

We now identify the branch coefficient.  Since
\(\widetilde s\in[0,p)\), equation (A.9) gives
\[
  \sg(\ell)=\sg(\widetilde s+p\ell)=\sg(V_w).
\tag{A.10}
\]
Indeed this is immediate for \(\ell\ge0\), while for \(\ell\le -1\) one has
\(\widetilde s+p\ell<0\).  Likewise, using (A.8),
\[
\begin{aligned}
  \Gamma_{\rho,\eta}
  &=p\Delta\eta+\Delta\rho+A_0-np \\
  &=\Delta(\rho+p\eta)+A_0-np \\
  &=\Delta(U_w-\alpha)+A_0-np \\
  &=\Delta U_w+\frac{pB}{2}.
\end{aligned}
\tag{A.11}
\]
If \(U_w\ge0\), then \(\Gamma_{\rho,\eta}>0\).  If \(U_w<0\), then
\(U_w=\widetilde r+p\eta\) with \(\eta\le -1\), hence
\(U_w\le \widetilde r-p\le -1/2\), and therefore
\[
  \Gamma_{\rho,\eta}\le -\frac{\Delta}{2}+\frac{pB}{2}
  =-\frac{np}{2}<0.
\]
Thus
\[
  \sg(\Gamma_{\rho,\eta})=\sg(U_w).
\tag{A.12}
\]
Combining (A.5), (A.10), and (A.12), the branch coefficient becomes
\[
  c_{\rho,\eta}(\ell)
  =\frac{\sg(V_w)+\sg(U_w)}2
  =\sgtwo(U_w,V_w).
\tag{A.13}
\]

It remains to compute the \(q\)-exponent.  The total exponent of the term
indexed by \((\rho,\sigma,\iota,\eta,\ell)\) is
\[
\begin{aligned}
\mathcal E
=&H(\rho,\sigma)
 +np^2C_2(\iota)+np(\sigma-\rho-n)\iota \\
&+(\Delta\sigma-A_0)\eta
 +(p\Delta\eta+\Delta\rho+A_0-np)\ell.
\end{aligned}
\tag{A.14}
\]
Using (A.6) and (A.7), a direct expansion gives
\[
  \mathcal E=H(R,S)+N C_2(w+1).
\tag{A.15}
\]
The following Sage snippet is included as an audit aid; the
identity is an elementary polynomial expansion:
\begin{verbatim}
var('rho sigma iota eta ell n p')
C2 = lambda z: z*(z-1)/2
B = n + p
Delta = p*(2*n + p)
N = n*Delta
A0 = n*p*(2*n + p + 1)/2
R = rho + n*iota + B*eta - n*ell
S = sigma - n*iota - n*eta + B*ell
w = ell - iota - eta
H = lambda X,Y: n*C2(X) + B*X*Y + n*C2(Y)
E = H(rho,sigma) + n*p^2*C2(iota) + n*p*(sigma-rho-n)*iota \
    + (Delta*sigma - A0)*eta \
    + (p*Delta*eta + Delta*rho + A0 - n*p)*ell
expand(E - H(R,S) - N*C2(w+1))
# output: 0
\end{verbatim}

Therefore, for fixed \(e,R,S\), equations (A.13)--(A.15) give
\[
\begin{aligned}
\relax[q^e(-x)^R(-y)^S]\,\theta^{\rm HM}_{n,p}
=\sum_{w\in\mathbb Z}
 \sgtwo\!\left(R+nw+\frac{n-1}{2},
               S-nw-\frac{n+1}{2}\right)
 \mathbf 1_{\,e=H(R,S)+N C_2(w+1)}.
\end{aligned}
\tag{A.16}
\]
Equivalently, since \((-x)^R(-y)^S=(-1)^{R+S}x^Ry^S\),
\[
\begin{aligned}
\relax[q^e x^R y^S]\,\theta^{\rm HM}_{n,p}
=&(-1)^{R+S}
\sum_{w\in\mathbb Z}
 \sgtwo\!\left(R+nw+\frac{n-1}{2},
               S-nw-\frac{n+1}{2}\right)  \\
&\qquad\qquad\times
 \mathbf 1_{\,e=H(R,S)+N C_2(w+1)}.
\end{aligned}
\tag{A.17}
\]
The sum is locally finite: for fixed \(e,R,S\), the equation
\(e=H(R,S)+N C_2(w+1)\) has at most two integer solutions \(w\).

But the coefficient of \(q^e x^R y^S\) in
\[
\bar\Theta_{n,p}
=\sum_{R,S,w\in\mathbb Z}(-x)^R(-y)^S
 \sgtwo\!\left(R+nw+\frac{n-1}{2},
               S-nw-\frac{n+1}{2}\right)
 q^{H(R,S)+N C_2(w+1)}
\]
is exactly the right-hand side of (A.17).  Hence
\[
  \theta^{\rm HM}_{n,p}(x,y,q)=\bar\Theta_{n,p}(x,y,q)
\]
coefficientwise.
